\chapter{보드를 검증한다}\label{ch:boardverify}

\ref{ch:verify}장의 원칙은 회로에도 그대로다. 믿을 수 있는 기준 모델을 두고, 같은 입력을 회로와 모델에 넣어
출력을 기계적으로, 대량으로 비교한다. 다른 것은 ``회로''가 HDL이 아니라 KiCad 넷리스트이고, 그것을 돌려 줄
시뮬레이터가 없다는 점이다. 그래서 이 프로젝트는 보드마다 작은 시뮬레이터를 직접 짰다. 이 장은 그 검증기들이
어떻게 생겼고 왜 그렇게 생겼는지를 본다. 시간 없는 접점 수준의 시뮬레이션, 기억이 있는 회로의 고정점,
보드 여러 장을 묶은 시험, 시간이 있는 사건 구동 시뮬레이션, 일부러 틀린 회로로 검증기를 시험하는 음성
시험, 그리고 KiCad 자신의 검사(ERC, DRC)다.

\section{검증기는 파일만 읽는다}\label{sec:v-indep}

\Why 생성기가 ``선 $i$의 NO를 $i + s$에 잇는다''고 썼다고 하자. 검증기가 생성기의 함수를 불러 ``선 $i$의 NO는
어디로 가는가''를 물으면 같은 답이 나온다. 생성기에 실수가 있으면 검증기도 같은 실수를 한다. 이런 검증은 아무것도
확인하지 않는다.

그래서 검증기는 생성기를 부르지 않는다. KiCad가 \emph{회로도 파일에서} 내보낸 넷리스트만 읽는다. 넷리스트에는
``어느 부품의 몇 번 핀이 어느 넷에 있는가''와 부품의 값·풋프린트만 있다. 검증기는 거기에 부품의 물리적 성질
(G6S-2의 핀 4는 극 1의 공통이고, 코일이 꺼지면 핀 3에 닿는다)과 기준 모델을 더해 회로를 직접 돌린다. 핀 기능도
생성기와 따로 적는다. 생성기가 핀 기능을 잘못 알고 있으면 둘이 어긋나 드러난다.

\General 소프트웨어에서 같은 사람이 같은 생각으로 짠 코드와 시험은 같은 오해를 공유한다. 사양에서 시험을 따로
쓰거나, 다른 사람이 쓰거나, 아예 다른 방법(예: 느리지만 단순한 구현)과 비교하는 것은 그 때문이다.

\Ours 검증기 머리말은 이 원칙을 적어 둔다. ``Model (netlist only, no knowledge of the generator)''(계수기),
``What is read from where (no knowledge of how the generator draws the tracks)''(오프셋 다이얼). 다이얼의 검증기는
동박의 모양을 \emph{기판 파일}에서 다각형으로 읽고, 빗의 위치는 기구 사양(\code{dialmech.py})에서, 기대값은
계수기 보드의 넷리스트를 돌려서 얻는다. 생성기와 검증기가 함께 읽는 것은 기준 모델과 기구 수치뿐이다.

\section{연결을 따라가기: 합집합--찾기}\label{sec:v-uf}

시간이 필요 없는 보드부터 보자. 시프터 보드는 제어값이 정해지면 릴레이들이 그 자리에 서 있고, 질문은 ``이
상태에서 입력 핀 $i$는 어느 출력 핀에 이어져 있는가''뿐이다.

\Why 이것은 그래프의 \term{연결 요소}{connected component}를 찾는 문제다. 넷 하나에 속한 핀들은 모두 한 덩어리다.
닫힌 접점 하나는 공통 핀과 NC(또는 NO) 핀의 두 덩어리를 합친다. 모든 넷과 닫힌 접점을 합치고 나면, 같은 덩어리에
든 핀끼리 전기적으로 이어진 것이다(그림~\ref{fig:v-uf}).

\begin{figure}[ht]
\centering
\begin{tikzpicture}[x=1cm,y=1cm,
    pin/.style={circle,draw,fill=white,inner sep=1.3pt},
    lab/.style={font=\tiny\ttfamily}]
  % nets as rounded boxes
  \foreach \n/\x/\y/\w/\t in {a/0/1.2/2.2/{J1.1, K101.4},b/3.8/2.2/3.0/{K101.3, K126.5, K201.4},c/3.8/0.2/3.0/{K101.5, K102.3, K202.4},
                               d/7.6/2.2/2.0/{K201.3, \ldots},e/7.6/0.6/2.0/{K201.5, \ldots}} {
    \node[draw=cgray,rounded corners=3pt,minimum width=\w cm,minimum height=7mm,font=\tiny\ttfamily,align=center] (\n) at (\x,\y) {\t};
  }
  \node[lab,above=0mm of a] {OIN\_A};
  \node[lab,above=0mm of b] {O1\_A};
  \node[lab,below=0mm of c] {O1\_B};
  \node[lab,above=0mm of d] {O2\_A};
  \node[lab,below=0mm of e] {O2\_C};
  \draw[line width=1.4pt,cblue] (a.east) -- node[above,font=\tiny,text=cblue] {COM--NC 닫힘} (b.west);
  \draw[dashed,cgray] (a.east) -- node[below,font=\tiny,text=cgray] {COM--NO 열림} (c.west);
  \draw[line width=1.4pt,cblue] (b.east) -- (d.west);
  \draw[dashed,cgray] (b.east) -- (e.west);
  \begin{pgfonlayer}{background}
    \fill[cbluebg,rounded corners=6pt] ($(a.south west)+(-0.15,-0.15)$) rectangle ($(d.north east)+(0.15,0.45)$);
  \end{pgfonlayer}
  \node[note,anchor=west,align=left] at (8.8,1.4) {파란 바탕: 한 덩어리\\= J1.1에서 갈 수 있는 곳};
\end{tikzpicture}
\caption{합집합--찾기의 생각. 상자 하나가 넷(같은 넷의 핀은 처음부터 한 덩어리), 선이 접점이다. 제어값이 0이면
단 1의 릴레이 K101이 놓여 있어 공통(핀 4)이 NC(핀 3)에 닿는다. 닫힌 접점만 덩어리를 합친다.}
\label{fig:v-uf}
\end{figure}

\General \term{합집합--찾기}{union--find}는 ``이 둘을 한 덩어리로 합쳐라''와 ``이것은 어느 덩어리인가''를 거의 상수
시간에 하는 자료 구조다. 덩어리마다 대표 하나를 두고, 원소는 부모를 따라 올라가면 대표에 닿는다. 지도의 섬을
세거나, 미로의 두 칸이 이어졌는지 보거나, 네트워크에서 끊긴 곳을 찾을 때 쓴다.

\Ours 시프터 검증기의 자료 구조는 열 줄이다. 따라 올라가는 길을 반씩 줄여(\code{p[a] = p[p[a]]}) 다음 찾기를
빠르게 한다.

\begin{lstlisting}[language=Python,caption={\code{verify\_shifter.py}: 합집합--찾기.},label={lst:v-uf}]
class UF:
    def __init__(self):
        self.p = {}

    def find(self, a):
        self.p.setdefault(a, a)
        while self.p[a] != a:
            self.p[a] = self.p[self.p[a]]
            a = self.p[a]
        return a

    def union(self, a, b):
        self.p[self.find(a)] = self.find(b)
\end{lstlisting}

제어값 26가지마다 넷을 합치고, 그 값에서 코일이 켜지는 릴레이를 정하고, 극마다 닫힌 쪽을 합친다. 켜지는 릴레이는
``코일의 핀 1이 제어 커넥터 J5의 해당 핀과 같은 넷인 릴레이''로 \emph{넷리스트에서} 찾는다.

\begin{lstlisting}[language=Python,caption={\code{verify\_shifter.py}: 제어값 \code{a}에서의 연결과 기대값.},label={lst:v-shifter}]
    for a in range(N):
        uf = UF()
        for nodes in nl.nets.values():
            for nd in nodes[1:]:
                uf.union(nodes[0], nd)
        energised = set()
        for b in range(BITS):
            if a >> b & 1:
                energised |= stage_coils[b]
        for r in relays:
            for com, nc, no in POLES:
                uf.union((r, com), (r, no if r in energised else nc))
        ...
        for i in range(N):
            tests = [
                ("outbound", "J1", "J3", shift_binary(i, a, False)),
                ("return", "J2", "J4", shift_binary(i, a, True)),
                ("outbound reversed", "J3", "J1", shift_binary(i, a, True)),
                ("return reversed", "J4", "J2", shift_binary(i, a, False)),
            ]
            for what, src, dst, want in tests:
                got = reach(cpin(src, i))
                if got != {cpin(dst, want)}:
                    ...
\end{lstlisting}

눈여겨볼 것은 비교가 \code{got != \{cpin(dst, want)\}}라는 점이다. ``기대한 출력 핀에 닿는가''가 아니라 ``닿는
신호 핀이 \emph{정확히} 그것 하나인가''를 본다. 다른 출력이나 다른 입력에도 닿으면 샛길이나 단락이고, 실패다.
\ref{sec:relay-logic}절에서 본 대로 접점은 양방향이라 이런 실수가 쉽게 생긴다. 기대값 \code{shift\_binary}는
\code{relay.rs}의 같은 이름 함수를 Python으로 옮긴 것이다. 네 방향(가는 길, 오는 길, 거꾸로 끼운 두 방향)을
26선 × 26값 모두 본다. 같은 검사가 제어 핀이 신호 핀에 닿지 않는지, 금속 틀(\code{CHASSIS})이 어떤 제어값에서도
신호에 닿지 않는지도 본다.

\section{기억이 있는 회로: 고정점까지 되풀이}\label{sec:v-fixpoint}

\Why 시프터에서는 릴레이의 상태가 바깥(제어값)에서 정해졌다. 계수기는 다르다. 릴레이의 코일이 다른 릴레이의
접점을 지나 전원을 받고, 자기 유지로 자기 접점에서 전원을 받기도 한다(\ref{sec:relay-latch}절). 릴레이의 상태가
연결을 정하고 연결이 다시 릴레이의 상태를 정한다. 이런 회로는 한 번에 풀리지 않는다. ``지금 상태에서 각 코일에
전원이 닿는가''를 계산하고, 그 결과대로 릴레이를 바꾸고, 다시 계산하기를 \emph{더 바뀌는 것이 없을 때까지}
되풀이한다. 더 바뀌지 않는 상태를 \term{고정점}{fixed point}이라 한다.

\General 표~\ref{tab:v-fix}는 \ref{sec:relay-latch}절의 시작/정지 회로를 이 방법으로 푼 것이다. 기억이 상태
벡터에 들어 있다는 것이 핵심이다. 시작 버튼을 놓은 뒤에도 $K$가 켜져 있는 것은, 그 직전 상태에서 $k$가 닫혀
있었기 때문이다. 같은 입력(버튼 둘 다 안 누름)이라도 이전 상태에 따라 고정점이 다르다.

\begin{table}[ht]
\centering
\caption{시작/정지 회로를 고정점 되풀이로 풀기. ``원함''은 그 상태에서 코일에 전원이 닿는가이다.}
\label{tab:v-fix}
\small
\begin{tabular}{llccl}
\toprule
입력 & 되풀이 & $K$ (지금) & $K$ (원함) & \\
\midrule
시작 누름 & 1 & 0 & 1 & 시작 버튼을 거쳐 닿음 $\to$ $K$ = 1 \\
          & 2 & 1 & 1 & 바뀌지 않음: 고정점 \\
버튼 놓음 & 1 & 1 & 1 & $k$를 거쳐 닿음: 고정점(기억) \\
정지 누름 & 1 & 1 & 0 & 정지 버튼이 끊음 $\to$ $K$ = 0 \\
          & 2 & 0 & 0 & 고정점 \\
버튼 놓음 & 1 & 0 & 0 & $k$가 열려 있어 닿지 않음: 고정점 \\
\bottomrule
\end{tabular}
\end{table}

\Ours 계수기 검증기의 \code{System}은 보드 여러 장을 한 회로로 묶어 이 일을 한다. 먼저 전원(+24~V와 바깥에서 거는
선)에서 출발해 닫힌 접점(양방향)과 다이오드(애노드에서 캐소드로만)를 따라 갈 수 있는 곳을 모두 찾는다. 0~V에
닿으면 단락이다.

\begin{lstlisting}[language=Python,caption={\code{verify\_counter.py}: 전원에서 갈 수 있는 곳과 고정점.},label={lst:v-settle}]
    def reach(self, drive_nets):
        ...
        while stack:
            n = stack.pop()
            for m, r, need in adj[n]:
                if seen[m]:
                    continue
                if r >= 0 and state[r] != need:
                    continue
                # r == -2: directed diode edge stored only at the anode, r == -1: fixed wire (unused)
                if m == gnd:
                    short = True
                    continue
                seen[m] = 1
                stack.append(m)
        return seen, short

    def settle(self, high: list[str], rng: random.Random | None = None, maxit: int = 400):
        """Apply the externally driven nets named in `high`, iterate to a fixed point. Returns the reach map."""
        dn = [n for h in high for n in self.drive[h]]
        state, coilp = self.state, self.coilp
        for _ in range(maxit):
            seen, short = self.reach(dn)
            if short:
                raise Fault("short circuit +24 V -> 0 V through contacts/diodes")
            want = [seen[c] for c in coilp]
            if want == state:
                return seen
            if rng is None:
                state[:] = want
            else:
                un = [i for i in range(len(state)) if want[i] != state[i]]
                i = rng.choice(un)
                state[i] = want[i]
        raise Fault("no fixed point (oscillation)")
\end{lstlisting}

접점의 간선에는 ``이 릴레이가 이 상태일 때만 지나갈 수 있다''는 조건(\code{need})이 붙어 있다. 다이오드는 애노드
쪽에만 간선을 저장해 한 방향으로만 따라간다. 400번을 되풀이해도 고정점이 없으면 회로가 \term{발진}{oscillation}하는
것이다. 릴레이가 자기 코일을 끊는 접점을 가지면 그렇게 된다.

\paragraph{두 가지 되풀이: 동기와 비동기.} \code{rng}가 없으면 ``바뀌어야 할 릴레이를 모두 한꺼번에'' 바꾼다.
\code{rng}가 있으면 바뀌어야 할 것 가운데 \emph{무작위로 하나만} 바꾸고 다시 계산한다. 실제 릴레이는 같은 순간에
움직이지 않으므로, 어떤 순서로 움직여도 같은 고정점에 닿아야 한다. 무작위 순서로 수만 번 돌려 결과가 늘 같으면
회로에 경쟁(\ref{sec:relay-counter}절)이 없다는 근거가 된다. 계수기는 두 방식 모두로 17,576클럭을 통과했다.

한 클럭은 이 고정점 계산 몇 번이다. 확정 상(\code{CK})을 걸고 풀고, 평상 상(\code{NCK})을 걸고 푼다. 그 사이마다
모든 유닛의 시프터 제어선과 자리마다 슬레이브·마스터의 값을 기준 모델과 비교한다.

\begin{lstlisting}[language=Python,caption={\code{verify\_counter.py}: 전체 구동. 클럭마다 두 시점에서 기준 모델과 비교.},label={lst:v-run}]
    refs = [Counters(k) for k in ks]
    seen = do_home(s, rng_mode)
    check_state(s, seen, refs, f"{label} after HOME")
    for step in range(1, clocks + 1):
        for c in refs:
            c.clock()
        seen1, seen2 = do_clock(s, rng_mode, skew)
        check_state(s, seen1, refs, f"{label} clock {step} (commit)", internal=False)
        check_state(s, seen2, refs, f"{label} clock {step} (release)")
\end{lstlisting}

\code{HOME}(프리셋) 시험은 릴레이 상태를 \emph{무작위로} 섞은 데서 시작한다. 전원을 넣은 직후 릴레이가 어떤
상태일지 모르므로, 어떤 상태에서 시작해도 \code{HOME}이 올바른 값을 만들어야 한다.

\section{보드를 묶어 기계 전체를}\label{sec:v-compose}

\Why 보드 한 장이 옳아도 보드 사이의 약속(커넥터의 핀 배치, 케이블의 방향, 카드를 꽂는 방향)이 틀리면 기계는
틀린다. 그래서 검증의 마지막 단계는 실제로 만들어질 보드들의 넷리스트를 케이블대로 이어 \emph{기계 전체}를 한
회로로 만들고, 그 출력을 기준 모델과 비교하는 것이다.

\General \ref{ch:verify}장의 층별 시험과 같다. 아래층(보드 한 장)이 맞는다는 것을 먼저 확인해야 위층(유닛 전체)이
틀렸을 때 원인을 좁힐 수 있다. 그리고 위층의 시험은 아래층 시험이 볼 수 없는 것, 곧 연결부를 본다.

\Ours 네 가지가 있다(그림~\ref{fig:v-compose}).

\begin{itemize}
\item \textbf{스크램블러 유닛 전체}(\code{wiring-cards/verify\_cards.py}). 시프터 보드 넷리스트 네 벌, 고른 로터 카드
세 장, 반사판 카드 한 장을 리본 케이블 14가닥대로 잇는다. 케이블 표는 카드가 꽂히는 자세에서 \emph{도출}하고,
그것이 조립 문서의 표와 같은지도 본다. 위치마다 \code{from\_cores}로 네 시프터의 제어값을 정하고, 입력 26선이 각각
정확히 출력 한 선에만 닿는지, 그 치환이 Rust 코어의 스크램블러(\ref{ch:rust-bombe}장)와 같은지 비교한다.
Enigma~I의 로터 순서 60개 × 17,576위치 = 1,054,560위치를 모두 돌려 보았다(약 13분).
\item \textbf{계수기 공통 보드 + 유닛 보드 셋}(\code{verify\_counter.py}). 오프셋 셋을 고른 유닛 셋을 공통 보드에
잇고 17,576클럭을 돌려, 아홉 번에 걸쳐 오프셋 26가지를 모두 지나간다.
\item \textbf{계수기 + 위치 표시}(\code{verify\_display.py}). 계수기 공통 보드의 넷리스트를 17,576클럭 돌리며 위치
커넥터를 표시 보드에 잇고, 켜진 램프 셋이 기준 모델의 $(l, m, r)$인지 본다.
\item \textbf{다이얼 동박 × 빗 × 계수기}(\code{verify\_dial.py}). 기판 파일에서 동박 다각형을 읽고, 빗 끝 11개를
기구 사양대로 놓아 어느 넷이 이어지는지 기하로 구한다. 그 연결을 계수기 유닛 보드의 넷리스트에 꽂아
\code{HOME}을 시뮬레이션하고, 프리셋이 기준 모델의 $(k, -k)$인지 본다. 멈춤 위치마다 회전 오차 $\pm 3°$, 반지름
오차 $\pm 0.3$~mm의 모서리 아홉 곳에서 같은 코드가 나와야 하고, 위치 사이를 훑는 동안 비트마다 이웃 두 코드
밖의 값이 나오면 안 된다.
\end{itemize}

\begin{figure}[ht]
\centering
\begin{tikzpicture}[x=1cm,y=1cm,
    bd/.style={blkn,minimum width=18mm,minimum height=8mm,font=\scriptsize},
    rf/.style={blkg,minimum width=26mm,minimum height=9mm,font=\scriptsize}]
  \node[bd] (s1) at (0,0) {시프터 ×4};
  \node[bd] (w) at (2.4,0) {로터 카드 ×3};
  \node[bd] (u) at (4.8,0) {반사판 카드};
  \draw[<->] (s1) -- (w); \draw[<->] (w) -- (u);
  \node[draw,dashed,cgray,fit=(s1)(w)(u),inner sep=3mm,label={[font=\scriptsize]above:넷리스트 + 리본 14가닥 = 유닛 하나}] (unit) {};
  \node[rf] (core) at (2.4,-2.0) {Rust 스크램블러};
  \node[blkr,minimum width=20mm,minimum height=9mm,font=\scriptsize] (cmp) at (6.6,-2.0) {치환 비교};
  \draw[->] (unit.south east) -- (cmp);
  \draw[->] (core) -- (cmp);
  \node[bd] (cc) at (9.0,0) {계수기 공통};
  \node[bd] (cu) at (11.4,0) {유닛 보드 ×3};
  \node[bd] (dd) at (9.0,-1.4) {위치 표시};
  \node[bd] (di) at (11.4,-1.4) {다이얼 동박};
  \draw[<->] (cc) -- (cu); \draw[->] (cc) -- (dd); \draw[->] (di) -- (cu);
  \node[rf] (cnt) at (10.2,-3.0) {\code{Counters} (기준)};
  \draw[->,dashed] (cnt) -- (dd); \draw[->,dashed] (cnt) -- (cu);
\end{tikzpicture}
\caption{보드를 묶어 시험하기. 왼쪽은 스크램블러 유닛, 오른쪽은 계수기와 그에 물리는 보드들이다. 모두 KiCad가
내보낸 실제 넷리스트(다이얼은 실제 기판 파일)로 묶는다.}
\label{fig:v-compose}
\end{figure}

\begin{lstlisting}[language=Python,caption={\code{verify\_cards.py}: 위치마다 유닛 전체의 치환을 기준과 비교.},label={lst:v-cards}]
    for l, m, r, k in positions:
        n += 1
        rk = (r + k) % N
        exp = oracle.scrambler(order, l, m, rk) if oracle else scrambler_port(order, l, m, rk)
        ...
        perm, problems = unit.permutation(from_cores(l, m, r, k))
        if problems or perm != exp:
            fails.append(...)
\end{lstlisting}

이 시험이 실제로 무엇을 잡았는지가 이 방법의 값어치를 보인다. 유닛의 맨 앞 시프터는 거꾸로 끼우므로 다른
카드와 반대쪽에서 보인다. 표준 기판을 그대로 쓰면 그 시프터에 꽂히는 리본 네 가닥의 핀 1이 반대 끝에 온다.
D-sub 커넥터는 모양이 비대칭이라 리본을 뒤집어 꽂을 수 없다. ``커넥터만 기판 뒷면에 납땜하면 되지 않나''라는 안은
\code{verify\_cards.py}가 기각했다. 뒷면 커넥터의 핀 배열을 풋프린트의 실제 좌표로 계산해 넣었더니 선 A--S가
거꾸로 이어지고 T--Z는 끊겨, 600가지 위치에서 한 번도 맞지 않았다. 그래서 커넥터를 거울 배치한 전용 기판
변형을 만들었다(\ref{sec:gen-place}절).

\section{시간이 있는 시뮬레이션}\label{sec:v-time}

계수기는 시간 없이 검증했다. ``각 상에서 고정점이 무엇인가''만 물었다. 그러나 순서기(\ref{sec:t-sequencer}절)와
테스트 레지스터(\ref{sec:t-register}절)는 시간이 본질이다. 질문이 ``\code{STROBE}는 몇 ms인가'', ``정지 릴레이가
\code{JUDGE} 안에 붙는가''이기 때문이다.

\Why 그래서 이 둘의 검증기는 \term{사건 구동}{event-driven} 시뮬레이션을 한다. 시간을 일정한 칸으로 나누지 않고,
``무언가가 바뀌는 시각''들의 목록(우선순위 큐)을 두고 가장 이른 사건부터 처리한다. 사건을 처리하면 회로를 다시
풀고(연결은 앞 절과 같은 방법), 그 결과로 생기는 새 사건(이 릴레이가 몇 ms 뒤에 붙는다)을 목록에 넣는다.

릴레이의 시간 모형은 결정 0001의 표(\ref{tab:t-relay})다.
\begin{itemize}
\item 코일에 전원이 닿으면 2.5--5~ms 뒤에 붙는다. 그 사이에 전원이 끊기면 \emph{붙지 않는다}. 짧은 펄스는 걸러진다.
\item 전원이 끊기면 1.5--4~ms(다이오드 + 제너) 또는 3--10~ms(다이오드만) 뒤에 놓인다.
\item 붙거나 놓일 때 접점은 0.2--0.6~ms 동안 공중에 떠 있다(먼저 끊고 나중에 잇기).
\end{itemize}
범위 안의 어디를 고를지는 실행 방식이 정한다. 무작위, 모두 가장 느리게, 모두 가장 빠르게, 그리고 \term{모서리}{corner}
(릴레이마다 최소나 최대 가운데 하나를 무작위로)다. 하한 규칙은 가장 나쁜 조합에서 깨지는데, 그 조합은 대개
``어떤 것은 가장 빠르고 어떤 것은 가장 느린'' 모서리에 있다.

\begin{lstlisting}[language=Python,caption={\code{seqsim.py}: 코일이 바뀌면 완료 사건을 예약하고, 완료 전에 되돌아가면 취소한다.},label={lst:v-relay}]
    def draw(self, lo, hi):
        if self.mode == "slow":
            return hi
        if self.mode == "fast":
            return lo
        if self.mode == "corner":
            return lo if self.rng.random() < 0.5 else hi
        return self.rng.uniform(lo, hi)

    def update_relays(self):
        for r in self.relays:
            on = self.lev.get(self.find(r.coil), 0) == 24
            if r.transit:
                continue
            if on != r.state:
                if r.pend is None or r.pend_to != on:
                    t = self.t + self.draw(*(r.op if on else r.rel))
                    r.pend = self.schedule(t, "relay", r)
                    r.pend_to = on
            else:
                r.pend = None          # coil returned to the old state before the move was completed

    def apply_relay(self, serial, r: Relay):
        if r.pend != serial:
            return
        r.pend = None
        r.transit = True
        self.schedule(self.t + self.rng.uniform(0.2, 0.6), "land", (r, r.pend_to))
\end{lstlisting}

예약한 사건에는 일련번호가 있다. 코일이 원래 상태로 돌아가면 \code{r.pend}를 지워 두고, 나중에 그 사건이 큐에서
나왔을 때 번호가 맞지 않으면 무시한다. 큐에서 사건을 꺼내 지우는 대신 ``이미 무효''로 표시하는 흔한 기법이다.
\code{transit} 동안은 그 릴레이의 공통이 어디에도 닿지 않는 것으로 푼다.

\paragraph{555도 넷리스트에서.} 순서기 시뮬레이터는 박자를 상수로 두지 않는다. 넷리스트에서 555의 방전 핀과 문턱
핀 사이 저항을 \emph{그 순간 닫혀 있는 접점으로 이어진 것까지} 더해 $R_2$를 구하고, 커패시터 전압을 지수 함수로
이어서 계산한다. 그래서 관찰 모드 저항을 단락하는 릴레이 접점이 언제 열리고 닫히는지가 그대로 박자에 반영된다
(\ref{sec:t-sequencer}절). 커패시터에는 $\pm$3\%의 공차를, 틱마다 $\pm$0.3\%의 흔들림을 넣는다.

\paragraph{규칙을 커넥터 핀에서 잰다.} 시뮬레이션이 끝나면 검증기는 커넥터 핀의 오르내림 기록에서 위치 주기마다
\code{STROBE}, \code{JUDGE}, \code{CLEAR}를 짝짓고 L1--L9를 하나씩 잰다.

\begin{lstlisting}[language=Python,caption={\code{verify\_sequencer.py}: 위치 주기마다 커넥터 면의 규칙.},label={lst:v-rules}]
L1, L3, L5 = 7.5, 7.5, 12.5
L2, L4, L6 = 11.5, 1.5, 1.5
...
    for k, ((s0, s1, _), (j0, j1, _), (c0, c1, _)) in enumerate(cyc):
        m = {"S": s1 - s0, "J": j1 - j0, "C": c1 - c0, "JS": j0 - s0, "CJ": c0 - j1}
        ...
        if m["S"] < L1:
            fails.append(f"L1: STROBE {m['S']:.2f} ms at t={s0:.1f}")
        if m["JS"] < L2:
            fails.append(f"L2: JUDGE rise only {m['JS']:.2f} after STROBE rise at t={s0:.1f}")
        ...
\end{lstlisting}

짝이 맞지 않는 것도 실패다. \code{STROBE} 뒤에 \code{JUDGE}와 \code{CLEAR}가 오지 않으면 L7, 앞에 \code{STROBE}가
없는 \code{JUDGE}는 L8 위반이다. 이 밖에 시험한 위치가 0, 1, 2, \ldots로 빠짐없이 이어지는지(정지와 재출발을
건너도 중복·건너뜀이 없는지), 놓친 정지나 거짓 정지가 없는지, \code{CK}와 \code{NCK}가 동시에 높지 않은지를 본다.
순서기는 42가지 시나리오(정지, 재출발, \code{HOME}, 대기, 비상 정지 등)와 조작 시나리오 40가지에 더해 17,576위치
전체를 한 번 돌렸다. 넷리스트 시뮬레이션으로 1,130초에 해당하는 운전을 89초에 계산한다.

\begin{lesson}{\code{HOME} 끝의 \code{CLEAR} 한 번}
비상 정지는 위치 주기의 어디서나 전원을 끊는다. \code{STROBE} 뒤, \code{CLEAR} 전에 끊기면 테스트 레지스터의 감지
릴레이가 붙은 채 남는다(감지 릴레이는 비상 정지로 끊기지 않는 전원에 있다). 그 뒤 \code{HOME}으로 처음 위치로
돌아가 출발하면, 첫 \code{STROBE}가 낡은 감지 릴레이 위에 덧붙는다. 남은 것이 26선 모두였다면 위치 0이 통과로
보여 그 위치의 정지를 놓친다. 순서기 시뮬레이션이 이 경우를 재현했고, 결정 0003은 ``\code{HOME}은 끝에 \code{CLEAR}를
정확히 한 번 낸다''는 규칙을 더했다. 지금은 이 \code{CLEAR}를 빼는 변형이 음성 시험에 들어 있어, 누가 지우면
``위치 0의 정지를 놓침''으로 걸린다. 기어식 제어 보드도 같은 규칙을 같은 방식으로 시험한다.
\end{lesson}

\section{일부러 틀린 회로: 음성 시험}\label{sec:v-mutants}

\Why 검증기가 ``통과''를 냈다고 하자. 그것은 회로가 옳다는 뜻일 수도 있고, 검증기가 아무것도 보지 않는다는 뜻일
수도 있다. 비교하는 줄에 오타가 있어 늘 같다고 판정하거나, 기대값을 회로에서 거꾸로 읽어 와 늘 맞거나,
시뮬레이터가 어떤 부품을 무시하고 있을 수 있다. 통과만 보아서는 둘을 구별할 수 없다.

그래서 \emph{일부러 틀린} 회로를 넣어 검증기가 실패를 내는지 본다. 이것을 \term{음성 시험}{negative test}이라 한다.
틀린 회로가 통과하면 검증기에 구멍이 있는 것이다. 어떤 실수를 넣을지는 ``실제로 일어날 법한 실수''에서 고른다.
핀 두 개를 바꿔 다는 것, NC와 NO를 바꾸는 것, 다이오드를 거꾸로 다는 것, 케이블을 바꿔 꽂는 것.

\General 소프트웨어에서는 이것을 \term{돌연변이 시험}{mutation testing}이라 한다. 코드의 \code{<}를 \code{<=}로,
\code{+}를 \code{-}로 바꾼 ``돌연변이''를 자동으로 만들어 시험을 돌리고, 시험이 잡지 못한 돌연변이가 있으면 시험이
그 부분을 보지 않는다는 뜻이다. 화재 경보기의 시험 단추도 같은 생각이다. 경보기가 조용한 것이 불이 없어서인지
고장이어서인지 알려면 가끔 일부러 울려 보아야 한다.

\Ours 모든 검증기가 마지막에 음성 시험을 돌리고, 하나라도 통과하면 전체를 실패로 끝낸다. 시프터 검증기의 변형은
넷리스트 \emph{텍스트}에서 핀 번호 두 개를 맞바꾸어 만든다.

\begin{lstlisting}[language=Python,caption={\code{verify\_shifter.py}: 음성 시험의 목록과 실행.},label={lst:v-neg}]
NEGATIVE = [
    ("swap two outbound OUT lines (J3 pins 3<->4)", "J3", "3", "4"),
    ("swap two outbound IN lines (J1 pins 1<->2)", "J1", "1", "2"),
    ("swap two return OUT lines (J4 pins 20<->25)", "J4", "20", "25"),
    ("swap NC/NO of stage-3 relay K305 pole 1 (pins 3<->5)", "K305", "3", "5"),
    ("swap control bits CTRL1<->CTRL2 (J5 pins 2<->3)", "J5", "2", "3"),
]
...
    for name, ref, pa, pb in NEGATIVE:
        bad = check(Netlist(swap_pins(text, ref, pa, pb)), quiet=True)
        print(f"  {'detected' if bad else 'MISSED  '}: {name} ({len(bad)} discrepancies)")
        ok &= bool(bad)
\end{lstlisting}

오프셋 다이얼의 음성 시험은 생성기의 숨은 키워드 인자(\ref{sec:gen-dial}절)로 \emph{틀린 다이얼을 새로 생성해서}
만든다. 비트 두 개의 링을 바꾼 코드, $a_2 = +k$로 잘못 계산한 코드, 글자 눈금을 한 칸 민 것, 호의 끝을 멈춤
쪽으로 0.42칸 옮긴 것, 다른 빗 간격에 맞춘 링, 끊어진 공통 링이다. 그 기판 파일을 정상과 똑같은 검증에 넣는다.

\begin{lstlisting}[language=Python,caption={\code{verify\_dial.py}: 틀린 다이얼을 생성해 검증에 넣는다.},label={lst:v-dial-neg}]
    variants = [
        ("code: a1 bits 0 and 1 on each other's ring", dict(code=swap_a1_01)),
        ("code: a2 = +k instead of -k", dict(code=lambda k: (k, k))),
        ("scale: letters shifted by one position", dict(letter_shift=1)),
        ("arc ends moved by 0.42 detent pitch (towards the detent)", dict(edge_shift=0.42 * M.ALPHA)),
        ("rings drawn for a comb of 1.5 mm pitch (tips 1.7 mm)", dict(pitch=1.5)),
        ("COM ring broken (20 deg gap at A)", dict(com_gap_deg=20.0)),
    ]
    for name, kw in variants:
        fpt = gen_dial.footprint_text(**kw)
        bt, _ = gen_dial_pcb.build(dial, fpt)
        _i, _j, errs = check_board(bt, dial, unit, work, "neg")
\end{lstlisting}

``호의 끝을 0.42칸 옮긴 것''은 0.5칸이 아니라 0.42칸이다. 끝이 멈춤 위치에서 0.08칸(약 1.1°) 떨어진 다이얼은
멈춤에서는 맞는 코드를 내지만 $\pm 3°$ 허용 오차 안에서 코드가 바뀐다. 검증기가 허용 오차의 모서리까지 실제로
보는지를 시험하는 변형이다. 이처럼 좋은 음성 시험은 ``아슬아슬하게 틀린'' 것이다. 크게 틀린 회로는 어떤
검증기로도 잡힌다.

표~\ref{tab:v-neg}는 보드별 음성 시험의 수다. 수가 많은 보드는 결정 0003(공통 전기 인터페이스)이 더한 항목(전원
입력의 보호 소자 방향, 신호 커넥터에 +24~V가 없음, 금속 틀과 0~V의 분리 등)도 하나씩 틀려 보기 때문이다.

\begin{table}[ht]
\centering
\caption{보드별 검증의 규모(각 README의 기록). 음성 시험은 모두 ``검출됨''이어야 통과한다.}
\label{tab:v-neg}
\small
\begin{tabularx}{\linewidth}{>{\raggedright\arraybackslash}p{30mm}>{\raggedright\arraybackslash}X>{\raggedright\arraybackslash}p{28mm}}
\toprule
보드 & 넷리스트 시뮬레이션 & 음성 시험 \\
\midrule
시프터 & 제어값 26 × 선 26 × 방향 4, 샛길·단락 없음 & 10 (+ 기판 2) \\
계수기 & 공통 + 유닛 셋, 17,576클럭, 오프셋 26가지, 동기·비동기 & 35 (+ 기판 3) \\
위치 표시 & 26$^3$ 위치, 무작위 릴레이 상태 3,000, 계수기와 끝에서 끝까지 & 16 (+ 기판 3) \\
배선 카드 & 유닛 전체 대 스크램블러, 최대 1,054,560위치 & 17 (+ 기판 3) \\
오프셋 다이얼 & 동박 기하 × 빗 × 계수기, 26위치 × 허용 오차 모서리 9 & 8 (+ 기판 3) \\
테스트 레지스터 & 사건 구동, 위치 주기 438, 정지 176, 최악·최선 시간 & 18 (+ 선 무전류 3) \\
순서기 & 사건 구동, 시나리오 82, 17,576위치 전체 & 31 \\
기어식 제어 보드 & 게이트·기구·펌웨어 모델·레지스터 연동 & 19 \\
\bottomrule
\end{tabularx}
\end{table}

\section{KiCad의 검사: ERC, DRC, 그리고 결정성}\label{sec:v-erc}

직접 짠 시뮬레이터 말고도 KiCad 자신의 검사를 명령행에서 돌린다.

\begin{itemize}
\item \textbf{ERC}(\term{전기 규칙 검사}{electrical rules check}). 회로도에서 ``이어지지 않은 핀'', ``출력끼리 이어진
핀'', ``전원을 주는 곳이 없는 전원 넷'' 같은 것을 찾는다. 일부러 쓰지 않는 핀에는 생성기가 \code{no\_connect} 표시를
달아 ``잊은 것이 아니라 비운 것''임을 밝힌다. 모든 보드가 모든 심각도에서 위반 0건이다.
\item \textbf{배치 대 회로도}. 기판 파일에서 패드별 넷을 내보내(IPC-D-356이라는 기판 시험용 넷리스트 형식, 또는
KiCad의 Python 모듈) 회로도의 넷리스트와 비교한다. 회로도의 부품이 모두 기판에 있고, 풋프린트가 같고, 모든 패드의
넷이 같아야 한다. 패드 두 개의 넷을 바꾼 기판, 풋프린트 하나를 뒷면으로 뒤집은 기판이 음성 시험이다.
\item \textbf{DRC}(\term{설계 규칙 검사}{design rule check}). 기판에서 부품의 차지 영역이 겹치는지, 구멍과 가장자리의
거리, 동박 사이의 간격을 본다. 배치만 한 기판이므로 ``이어지지 않은 항목''이 남는데, 그 수가 회로도에서 계산한
값, 곧 넷마다 (패드 수 $-$ 1)의 합과 같아야 한다. 패드 $n$개를 잇는 데는 선이 최소 $n-1$개 필요하기 때문이다.
시프터 보드는 1,125개다.
\item \textbf{결정성}. 생성기를 임시 폴더에 다시 돌려 저장소의 파일과 바이트 단위로 비교한다(\ref{sec:gen-uuid}절).
\end{itemize}

\section{시뮬레이션이 보지 못하는 것}

이 모든 검증은 \emph{설계}의 검증이다. 각 README는 무엇을 모형에 넣지 않았는지를 적어 둔다.

\begin{itemize}
\item 접점의 튐은 따로 모형에 없다. 붙는 시간에 안정 시간을 포함시켰을 뿐이다.
\item 코일 전압의 강하(배선과 다이오드를 지나며 줄어드는 전압)는 계산으로만 보았다. 시뮬레이터에서 코일은 켜지거나
꺼질 뿐이다.
\item 릴레이 시간 자체가 데이터시트와 추정에 기댄다(표~\ref{tab:t-relay}). 다이오드만 단 코일의 복귀 시간은 추정이다.
\item 반도체(MOSFET의 문턱, LED 전류)는 저항값으로 계산해 확인했고 아날로그 시뮬레이션은 하지 않았다.
\end{itemize}

그래서 각 설계 문서는 시제품에서 잴 것을 목록으로 남긴다. 감지 릴레이가 붙어 유지에 들기까지의 시간 분포,
놓이는 시간, 순서기 출력 릴레이의 실제 지연, 두 클럭 버스의 어긋남 같은 것이다. 시뮬레이션은 ``이 범위 안이면
옳다''를 보이고, 실측은 ``실제 부품이 그 범위 안에 있다''를 보인다. 둘이 함께 있어야 기계가 옳다.

\section{정리}

\begin{itemize}
\item 검증기는 생성기를 부르지 않고 KiCad가 내보낸 파일만 읽는다. 같은 실수를 공유하지 않기 위해서다.
\item 시간 없는 보드는 합집합--찾기로 연결을 구하고, 기대한 연결 \emph{하나만} 있는지 본다.
\item 기억이 있는 회로는 고정점까지 되풀이한다. 무작위 순서로 되풀이해도 같은 결과면 경쟁이 없다.
\item 보드를 케이블대로 묶어 기계 전체를 기준 모델과 비교한다. 연결부의 실수는 여기서만 잡힌다.
\item 시간이 본질인 보드는 사건 구동으로 돌리고, 릴레이 시간을 범위의 모서리까지 흔들며 하한 규칙을 커넥터 핀에서
잰다.
\item 일부러 틀린 회로가 모두 실패해야 검증기를 믿는다. 좋은 음성 시험은 아슬아슬하게 틀린 것이다.
\end{itemize}

\begin{keybox}[title={H부에서 쓴 방법들}]
\begin{itemize}
\item 접점은 양방향 길이다. 논리는 ``전원에서 닫힌 접점만 밟아 올 수 있는가''이고, 샛길을 늘 의심한다.
\item 기억은 자기 유지, 정확히 하나는 전환 접점, 경쟁 없는 계수는 마스터--슬레이브와 두 상의 클럭.
\item 흩어지는 부품에는 하한으로 규칙을 정하고, 느려져도 안전하게 만든다. 회로의 면과 커넥터의 면을 나눠 계산한다.
\item 보드는 기준 모델에서 생성하고, 생성과 무관한 검증기가 파일만 읽어 기준 모델과 대조한다. 그 검증기는 일부러
틀린 보드로 다시 시험한다.
\item 전원과 안전은 소프트웨어가 아니라 레일·퓨즈·접점의 사슬로 지킨다(다음 장, \ref{ch:44-power-safety}장).
\end{itemize}
\end{keybox}
